% Owner: TV3 (P5.3b) -- Reviewer: TV1
% The assignment limits the abstract to 200 words: problem, method, key numerical results, conclusion.
\begin{abstract}
We compare Minimax, alpha-beta pruning, and Monte Carlo tree search with UCT (MCTS-UCT) on $3\times 3$ Tic-Tac-Toe and on Gomoku (five in a row), with an extension to Xiangqi. All matches are run by one seeded runner, with the first move alternated and 95\% Wilson confidence intervals. On $3\times 3$, alpha-beta returns the Minimax value while visiting 96.7--98.7\% fewer nodes, depending on move ordering. MCTS-UCT never loses to any opponent strategy over 20 seeds from $I^{*} = 2\,000$ iterations per move, although it already draws every game against perfect Minimax at 1\,000. The exploration constant $c$ matters only for small budgets: on $9\times 9$ Gomoku with 100 iterations, the win rate falls from 96.5\% to 88.5\% as $c$ grows from 0.5 to 2. Exploration bonuses that decay as $1/N$ weaken MCTS, while a heuristic prior raises its win rate from 49\% to 65.5\%. Informed rollouts win 65.5--90\% against random ones at equal iterations, but only a truncated rollout with an evaluation function still wins at equal time. On $15\times 15$ Gomoku, alpha-beta with a pattern evaluation beats MCTS in 594 of 600 games. Domain knowledge helps search only when it is cheap compared with the simulations it replaces.
\end{abstract}
